%%
%% 2026年度第2回学生研究発表会　Sticky Manifold PDMP 暫定案
%% UTF-8。platex sticky_pdmp_2027.tex を2回実行し，dvipdfmx -p a4 sticky_pdmp_2027.dvi でPDFを作成する。
%%

\documentclass[leqno,a4j]{jarticle}
%以下の２行は，式番号を章ごとに1.１，1.2，2.1，2.2，...とつけるための設定
\usepackage{amsmath}
\numberwithin{equation}{section}

\pagestyle{empty}

\begin{document}
\normalsize

\begin{center}

\textbf{\huge{ベイズモデル平均のための\\Sticky Manifold PDMP}}
\vspace{0.5cm}

%%% コース/専攻名・学年を以下に入力
%%% 2022年10月までに入学された方は「統計科学専攻　博士課程○年」（５年一貫制による学年（３～５年）を記入）として下さい．
%%% 2023年4月以降に博士後期課程に入学された方は「統計科学コース　博士後期課程○年」（５年一貫制による学年（３～５年）を記入）として下さい．
\large{統計科学コース　５年一貫制博士課程４年\\

%%% 発表者氏名
司馬　博文
\vspace{0.2cm}\\ 
%%% 主任指導教員名
\normalsize{主任指導教員：　鎌谷　研吾}

%%% 主任指導教員以外の共同研究者については，本文中に明記して下さい 
}

\end{center}


\section{モデルの次元も推定する}
一つの量子ビットの状態は，単位球内のベクトル $r$ で表せる．
球面 $\|r\|=1$ は純粋状態，球内部 $\|r\|<1$ は混合状態に対応する．
観測後に「純粋状態である確率」を知りたければ，球面に正の事前確率を置く必要がある．
球内の体積に関して密度をもつ事前分布を選ぶと，どれほど球面の近くに事後分布が集中しても，
球面そのものの事後確率はゼロのままである．

一般に，$d$ 次元領域 $D$ のモデルと，余次元 1 の滑らかな部分多様体 $M$ 上のモデルを考える．
両モデルが同じ位置 $x$ で尤度 $L(y\mid x)$ を共有するとき，混合次元の事後分布は
\begin{equation}
 \pi(dx\mid y)\ \propto\ L(y\mid x)
 \bigl\{w\,q_D(x)\,dx+(1-w)\,q_M(x)\,dS_M(x)\bigr\}
 \label{eq:posterior}
\end{equation}
となる．$w$ は $D$ の事前モデル確率，$q_D,q_M$ はそれぞれ体積と表面積に関する事前密度である．
計算上の課題は，次元の異なる両モデルを行き来し，それぞれの事後確率に見合う時間を過ごす軌道を作ることである．
本研究は Luke Hardcastle 氏（ケンブリッジ大学）および鎌谷研吾氏との共同研究である．

\section{境界で速度を切り替える}
Sticky Manifold PDMP は，モデルごとに区分的決定論的マルコフ過程（PDMP）を走らせ，
境界で二つの過程を接続する．PDMP はランダムなイベントの間，決定論的な軌道をたどる．
$M$ に到達すると，速度を接成分と法線成分に分け，法線方向の速さ $a=|v^\perp|$ を保持する．
位置は接方向に進み続ける．離脱時計が鳴ると，その時点の位置から法線方向の運動を再開する．

$M$ の両側から到達できる場合，保持した速さ $a$ と現在の位置 $x$ から離脱率を定める．
\begin{equation}
 \lambda(x,a)=a\,\frac{w\,q_D(x)}{(1-w)\,q_M(x)},\qquad x\in M .
 \label{eq:rate}
\end{equation}
法線方向の速さが大きい軌道ほど境界に多く流入するので，その分だけ早く離脱させる．
同じ $x$ で両モデルの密度を比べるため尤度が相殺され，モデル切替えには追加の尤度評価も採否判定も要らない．

\textbf{不変性定理．}
各モデル内の過程は，条件付き事後分布に速度を付加した所定の同時分布を保つとする．
付着・離脱時の速度更新は，位置と法線方向の速さを保存し，
入射・滞在・出射の速度分布を順に対応させるものとする．
所定の正則性条件の下で，接続した過程は式~\eqref{eq:posterior} を位置周辺分布とする不変分布をもつ
（Shiba et al. (2026)，定理 3.5）．
証明では，決定論的な流れに沿う Gauss--Green 公式を用い，$M$ への確率流入が
式~\eqref{eq:rate} による流出と釣り合うことを示す．

\newpage
\section{量子状態のモデル確率とモデル平均}
球内部と球面にそれぞれ事前確率 $1/2$ を置き，各モデル内では一様分布とする．
単位方向 $u$ に沿った測定値 $Y\in\{-1,1\}$ の確率は
$\Pr(Y=y\mid r,u)=(1+yu^{\mathsf T}r)/2$ である．
推定対象は純粋状態の事後確率 $p=\Pr(\|r\|=1\mid y)$ と，
モデル平均した状態 $m$ である．
\begin{equation}
 m=\mathrm{E}[r\mid y]
   =p\,\mathrm{E}[r\mid y,\|r\|=1]
   +(1-p)\,\mathrm{E}[r\mid y,\|r\|<1] .
 \label{eq:bma}
\end{equation}
エルゴード性の下で，$p$ は球面にいた時間の割合，
$m$ は球面と球内部を通じた $r$ の時間平均で推定できる．

この例では球面に内側からのみ到達するので，対称な速度分布では式~\eqref{eq:rate} の離脱率を半分にする．
体積密度 $3/(4\pi)$ と面積密度 $1/(4\pi)$ を代入すると $\lambda=3a/2$ となる．
標準ガウス速度の法線成分を積分し，定数になった離脱率の下で指数滞在時間をその平均に置き換えると，
球面にいる時間を毎回 $\sqrt{2\pi}/3$ に固定しても位置の事後分布が保たれる
（同論文，系 3.8--3.10）．一回の滞在時間は同じでも，球面に戻る頻度には観測データが反映される．

\textbf{数値実験．}
$r_\alpha=\alpha(1,1,1)^{\mathsf T}/\sqrt 3$ から，無作為な方向で独立に100回測定した．
各反復で同じ事後標本を初期値とし，各手法を1 CPU秒ずつ動かした．100反復から二乗平均平方根誤差（RMSE）を求めた．
RB-SM-FECMC は，球内では直線，球面では大円に沿って進む forward event-chain Monte Carlo に，上記の Rao--Blackwell 化（RB）を適用した手法である．
比較にはモデル間を可逆ジャンプ（RJ）で移動し，モデル内は同じ FECMC または
ハミルトニアン・モンテカルロ（HMC）で探索する手法を用いた．

\begin{center}
表1　純粋状態の事後確率 $p$ とモデル平均 $m$ の RMSE（単位 $10^{-4}$）\\[3pt]
\begin{tabular}{lrrrr}
\hline
 & \multicolumn{2}{c}{$\alpha=1.00$（純粋）} & \multicolumn{2}{c}{$\alpha=0.99$（混合）} \\
\cline{2-3}\cline{4-5}
手法 & $p$ & $m$ & $p$ & $m$ \\
\hline
RB-SM-FECMC（提案法）  & \textbf{4.91} & 3.30 & \textbf{4.87} & 3.48 \\
RJ-FECMC              & 9.90 & 6.28 & 15.9 & 5.37 \\
RJ-HMC（$\epsilon=0.08$） & 6.40 & 2.46 & 7.91 & 2.39 \\
RJ-HMC（$\epsilon=0.09$） & 6.25 & \textbf{2.25} & 8.32 & \textbf{2.36} \\
\hline
\end{tabular}
\end{center}
モデル間の移動回数をほぼ同じにしても，両データで $p$ の誤差は提案法が最小だった．
一方，$m$ の誤差は RJ-HMC が小さい．式~\eqref{eq:bma} に現れる各モデル内の条件付き平均も
RJ-HMC が精度よく推定しており，モデル確率とモデル平均で手法の順位が変わった
（同論文，表 2(a)）．
提案法は軌道生成中に対数尤度の値を評価せず，その微分を用いる．値と微分の評価を合計 $10^6$ 回に合わせると，両データで $p,m$ ともに提案法の誤差が最小となった（同表 (b)）．

\textbf{他の応用．}
20次元のベイズ主成分分析では，固有値と固有ベクトルを座標に取り，最小の正の固有値がゼロになる面を，隣接ランクを結ぶ余次元 1 の境界として表した（blow-up）．
信号共分散行列 $G$ の推定では，6設定すべてで $\mathrm{E}[G\mid y]$ の RMSE は
二つの FECMC 法のいずれかが最小だった．生存時間の外挿では，一般化ガンマ族に埋め込んだ
対数正規・ワイブル等の分布族をまたぎ，30年間の制限平均生存時間をモデル平均した．

\begin{center}
\textbf{\large{参考文献}}
\end{center}
\noindent
Shiba, H., Hardcastle, L. and Kamatani, K. (2026).
Sticky manifold piecewise deterministic Monte Carlo for Bayesian model averaging.
arXiv:2611.17008.

\end{document}
